% Learning objectives
\begin{frame}{Learning Objectives}
By the end of this chapter, you will understand:
\begin{itemize}
\item Performance optimization principles in digital logic design
\item Chisel control constructs: \texttt{when}, \texttt{MuxCase}, and \texttt{switch}
\item Synthesis implications of different coding styles
\item Hardware structures: multiplexers vs. priority logic
\item Timing, area, and power trade-offs
\item Best practices for performance-optimized design
\item Tool-specific synthesis considerations
\end{itemize}
\end{frame}

% Section 1: Introduction to Performance Optimization
\section{Introduction to Performance Optimization}

\begin{frame}{Performance Optimization in Digital Design}
\begin{columns}
\begin{column}{0.6\textwidth}
\concept{Optimization Goals}:
\begin{itemize}
\item \textbf{Timing Performance}: Minimize critical path delay
\item \textbf{Area Efficiency}: Reduce hardware resource usage
\item \textbf{Power Consumption}: Optimize for low power operation
\item \textbf{Synthesis Quality}: Generate efficient hardware structures
\end{itemize}

\vspace{0.3cm}
\concept{Design Considerations}:
\begin{itemize}
\item Choice of HDL constructs affects synthesis results
\item Different coding styles produce different hardware
\item Tool-specific optimizations and limitations
\item Trade-offs between competing objectives
\end{itemize}

\vspace{0.3cm}
\concept{Chisel-Specific Factors}:
\begin{itemize}
\item High-level constructs map to hardware structures
\item Functional programming paradigms influence synthesis
\item Type system provides optimization opportunities
\item Generator-based design enables parameterization
\end{itemize}
\end{column}
\begin{column}{0.4\textwidth>
\begin{center}
\textbf{Optimization Triangle}
\vspace{0.3cm}

\tikz{
% Triangle
\draw[thick] (0,0) -- (2,3) -- (4,0) -- cycle;

% Vertices
\node at (0,-0.3) {\textbf{Area}};
\node at (4,-0.3) {\textbf{Power}};
\node at (2,3.3) {\textbf{Performance}};

% Center
\node at (2,1) {Design};
\node at (2,0.7) {Space};

% Trade-off arrows
\draw[<->] (0.5,0.5) -- (3.5,0.5);
\draw[<->] (1,2.5) -- (3,2.5);
\draw[<->] (0.8,0.3) -- (1.8,2.7);
}
\end{center}
\end{column}
\end{columns}
\end{frame}

\begin{frame}{HDL Construct Impact on Hardware}
\begin{columns}
\begin{column}{0.6\textwidth}
\concept{Behavioral vs. Structural Implications}:
\begin{itemize}
\item \textbf{Behavioral}: High-level description of functionality
\item \textbf{Structural}: Direct mapping to hardware components
\item \textbf{Synthesis}: Translation from behavioral to structural
\item \textbf{Optimization}: Tool-dependent improvements
\end{itemize}

\vspace{0.3cm}
\concept{Control Structure Categories}:
\begin{itemize}
\item \textbf{Sequential Logic}: \texttt{when-elsewhen-otherwise}
\item \textbf{Parallel Logic}: \texttt{MuxCase} and multiplexers
\item \textbf{Hybrid Logic}: \texttt{switch} statements
\item \textbf{Conditional Logic}: Priority-based selection
\end{itemize}

\vspace{0.3cm}
\concept{Synthesis Considerations}:
\begin{itemize}
\item Tool interpretation of HDL constructs
\item Optimization passes and transformations
\item Target technology constraints
\item Design rule compliance
\end{itemize}
\end{column}
\begin{column}{0.4\textwidth>
\begin{center}
\textbf{HDL to Hardware Flow}
\vspace{0.3cm}

\tikz{
% Flow diagram
\node[draw, rectangle, minimum width=2cm, minimum height=0.6cm] (hdl) at (0,3) {HDL Code};
\node[draw, rectangle, minimum width=2cm, minimum height=0.6cm] (syn) at (0,2.2) {Synthesis};
\node[draw, rectangle, minimum width=2cm, minimum height=0.6cm] (opt) at (0,1.4) {Optimization};
\node[draw, rectangle, minimum width=2cm, minimum height=0.6cm] (hw) at (0,0.6) {Hardware};

\draw[->] (hdl) -- (syn);
\draw[->] (syn) -- (opt);
\draw[->] (opt) -- (hw);

% Side annotations
\node[right] at (2.2,3) {Behavioral};
\node[right] at (2.2,2.2) {Translation};
\node[right] at (2.2,1.4) {Improvement};
\node[right] at (2.2,0.6) {Structural};
}
\end{center}
\end{column}
\end{columns}
\end{frame}

% Section 2: Chisel Control Constructs
\section{Chisel Control Constructs}

\begin{frame}[fragile]{The \texttt{when} Construct}
\begin{columns}
\begin{column}{0.5\textwidth}
\concept{Syntax and Semantics}:
\begin{lstlisting}[style=chisel]
when(condition1) {
  // Execute when condition1 true
}.elsewhen(condition2) {
  // Execute when condition2 true
}.otherwise {
  // Default case
}
\end{lstlisting}

\vspace{0.3cm}
\concept{Characteristics}:
\begin{itemize}
\item Sequential evaluation of conditions
\item Priority-based logic structure
\item Natural for software-like if-else logic
\item Cascaded conditional implementation
\end{itemize}
\end{column}
\begin{column}{0.5\textwidth>
\concept{Hardware Generation}:
\begin{itemize}
\item Creates priority logic chain
\item Sequential condition checking
\item Early termination when condition met
\item Propagation delay through cascade
\end{itemize}

\vspace{0.3cm}
\concept{Use Cases}:
\begin{itemize}
\item Natural priority ordering
\item Sequential decision making
\item Default/fallback cases
\item Human-readable conditions
\end{itemize}

\vspace{0.3cm}
\concept{Performance Impact}:
\begin{itemize}
\item Longer critical path for later conditions
\item Compact area for early termination
\item Power efficient for sparse conditions
\end{itemize}
\end{column}
\end{columns}
\end{frame}

\begin{frame}[fragile]{The \texttt{MuxCase} Construct}
\begin{columns}
\begin{column}{0.5\textwidth}
\concept{Syntax and Semantics}:
\begin{lstlisting}[style=chisel]
val result = MuxCase(default, Seq(
  (condition1 -> value1),
  (condition2 -> value2),
  (condition3 -> value3)
))
\end{lstlisting}

\vspace{0.3cm}
\concept{Characteristics}:
\begin{itemize}
\item Parallel evaluation of conditions
\item Priority-free selection logic
\item Direct multiplexer mapping
\item Simultaneous condition checking
\end{itemize}
\end{column}
\begin{column}{0.5\textwidth>
\concept{Hardware Generation}:
\begin{itemize}
\item Creates multiplexer tree structure
\item Parallel logic evaluation
\item Fixed propagation delay
\item Predictable timing characteristics
\end{itemize}

\vspace{0.3cm}
\concept{Use Cases}:
\begin{itemize}
\item Mutually exclusive conditions
\item Performance-critical paths
\item Parallel decision making
\item Fixed-latency requirements
\end{itemize}

\vspace{0.3cm}
\concept{Performance Impact}:
\begin{itemize}
\item Consistent timing across conditions
\item Larger area for parallel evaluation
\item Higher power for all conditions
\end{itemize}
\end{column}
\end{columns}
\end{frame}

\begin{frame}[fragile]{The \texttt{switch} Construct}
\begin{columns}
\begin{column}{0.5\textwidth}
\concept{Syntax and Semantics}:
\begin{lstlisting}[style=chisel]
switch(selector) {
  is("b00".U) { result := a }
  is("b01".U) { result := b }
  is("b10".U) { result := c }
  is("b11".U) { result := d }
}
\end{lstlisting}

\vspace{0.3cm}
\concept{Characteristics}:
\begin{itemize}
\item Multi-way branching based on value
\item Flexible synthesis options
\item Can generate mux or priority logic
\item Tool-dependent optimization
\end{itemize}
\end{column}
\begin{column}{0.5\textwidth>
\concept{Hardware Generation Options}:
\begin{itemize}
\item \textbf{Multiplexer Tree}: For mutually exclusive cases
\item \textbf{Priority Logic}: For overlapping conditions
\item \textbf{Decoder Logic}: For wide selectors
\item \textbf{Hybrid Structures}: Tool-optimized combinations
\end{itemize}

\vspace{0.3cm}
\concept{Synthesis Factors}:
\begin{itemize}
\item Condition exclusivity
\item Selector width and encoding
\item Tool optimization capabilities
\item Target technology characteristics
\end{itemize}
\end{column}
\end{columns}
\end{frame}

% Section 3: Synthesis Implications
\section{Synthesis Implications}

\begin{frame}{Timing and Performance Analysis}
\begin{columns}
\begin{column}{0.6\textwidth}
\concept{Critical Path Considerations}:
\begin{itemize}
\item \textbf{\texttt{when}}: Sequential evaluation creates cascaded delays
\item \textbf{\texttt{MuxCase}}: Parallel evaluation with fixed delay
\item \textbf{\texttt{switch}}: Variable delay based on synthesis choice
\end{itemize}

\vspace{0.3cm}
\concept{Propagation Delay Analysis}:
\begin{itemize}
\item \textbf{Cascaded Logic}: Delay accumulates through chain
\item \textbf{Parallel Logic}: Single-level delay for all paths
\item \textbf{Multiplexer Tree}: Logarithmic delay with input count
\item \textbf{Priority Encoder}: Linear delay with condition count
\end{itemize}

\vspace{0.3cm}
\concept{Performance Optimization}:
\begin{itemize}
\item Choose construct based on timing requirements
\item Consider condition ordering for \texttt{when}
\item Use \texttt{MuxCase} for critical timing paths
\item Balance area and timing trade-offs
\end{itemize}
\end{column}
\begin{column}{0.4\textwidth>
\begin{center}
\textbf{Timing Comparison}
\vspace{0.3cm}

\tikz{
% Timing bars
\node at (1.5,3.5) {\textbf{Propagation Delay}};

% when construct
\draw[thick, fill=red!30] (0,3) rectangle (2.5,3.2);
\node at (-0.5,3.1) {\texttt{when}};
\node at (3,3.1) {Variable};

% MuxCase construct  
\draw[thick, fill=green!30] (0,2.6) rectangle (1.5,2.8);
\node at (-0.5,2.7) {\texttt{MuxCase}};
\node at (3,2.7) {Fixed};

% switch construct
\draw[thick, fill=blue!30] (0,2.2) rectangle (2,2.4);
\node at (-0.5,2.3) {\texttt{switch}};
\node at (3,2.3) {Tool-dep};

% Scale
\node at (0,1.8) {Fast};
\node at (2.5,1.8) {Slow};
\draw[->] (0,1.9) -- (2.5,1.9);
}
\end{center}
\end{column}
\end{columns}
\end{frame}

\begin{frame}{Area and Resource Utilization}
\begin{columns}
\begin{column}{0.6\textwidth}
\concept{Hardware Resource Requirements}:
\begin{itemize}
\item \textbf{\texttt{when}}: Compact for early termination cases
\item \textbf{\texttt{MuxCase}}: Larger due to parallel evaluation
\item \textbf{\texttt{switch}}: Variable based on synthesis strategy
\end{itemize}

\vspace{0.3cm}
\concept{Logic Utilization Patterns}:
\begin{itemize}
\item \textbf{Priority Logic}: Sequential gates, pruning opportunities
\item \textbf{Multiplexer Logic}: Parallel gates, full evaluation
\item \textbf{Decoder Logic}: One-hot encoding, distributed control
\end{itemize}

\vspace{0.3cm}
\concept{Optimization Strategies}:
\begin{itemize}
\item Order conditions by probability in \texttt{when}
\item Use \texttt{MuxCase} for mutually exclusive cases
\item Consider condition complexity and frequency
\item Leverage synthesis tool optimizations
\end{itemize}
\end{column}
\begin{column}{0.4\textwidth>
\begin{center}
\textbf{Area Comparison}
\vspace{0.3cm}

\tikz{
% Area bars
\node at (1.5,3.5) {\textbf{Hardware Area}};

% when construct
\draw[thick, fill=green!30] (0,3) rectangle (1.2,3.2);
\node at (-0.5,3.1) {\texttt{when}};
\node at (2.5,3.1) {Compact};

% MuxCase construct  
\draw[thick, fill=red!30] (0,2.6) rectangle (2.2,2.8);
\node at (-0.5,2.7) {\texttt{MuxCase}};
\node at (2.5,2.7) {Larger};

% switch construct
\draw[thick, fill=blue!30] (0,2.2) rectangle (1.7,2.4);
\node at (-0.5,2.3) {\texttt{switch}};
\node at (2.5,2.3) {Variable};

% Scale
\node at (0,1.8) {Small};
\node at (2.2,1.8) {Large};
\draw[->] (0,1.9) -- (2.2,1.9);
}
\end{center}
\end{column}
\end{columns}
\end{frame}

\begin{frame}{Power Consumption Analysis}
\begin{columns}
\begin{column}{0.6\textwidth}
\concept{Power Consumption Factors}:
\begin{itemize}
\item \textbf{Dynamic Power}: Switching activity and capacitance
\item \textbf{Static Power}: Leakage current in inactive logic
\item \textbf{Activity Factor}: Percentage of logic switching
\item \textbf{Clock Power}: Distribution and sequential elements
\end{itemize}

\vspace{0.3cm}
\concept{Construct-Specific Power Characteristics}:
\begin{itemize}
\item \textbf{\texttt{when}}: Lower power for sparse conditions
\item \textbf{\texttt{MuxCase}}: Higher power due to parallel evaluation
\item \textbf{\texttt{switch}}: Power depends on synthesis choice
\end{itemize}

\vspace{0.3cm}
\concept{Power Optimization Techniques}:
\begin{itemize}
\item Clock gating for inactive logic
\item Condition ordering for early termination
\item Power-aware synthesis settings
\item Activity-driven optimization
\end{itemize}
\end{column}
\begin{column}{0.4\textwidth>
\begin{center}
\textbf{Power Characteristics}
\vspace{0.3cm}

\tikz{
% Power components
\node at (1.5,3.5) {\textbf{Power Breakdown}};

% Dynamic power
\draw[thick, fill=red!30] (0,3) rectangle (2,3.2);
\node at (1,3.1) {Dynamic};

% Static power
\draw[thick, fill=blue!30] (0,2.7) rectangle (1.2,2.9);
\node at (0.6,2.8) {Static};

% Clock power
\draw[thick, fill=green!30] (0,2.4) rectangle (1.5,2.6);
\node at (0.75,2.5) {Clock};

% Activity factor impact
\node at (1.5,2) {\textbf{Activity Impact}};
\node at (1.5,1.7) {High activity → High power};
\node at (1.5,1.4) {Sparse conditions → Low power};
}
\end{center}
\end{column}
\end{columns}
\end{frame}

% Section 4: Hardware Structure Analysis
\section{Hardware Structure Analysis}

\begin{frame}{Multiplexer Tree Structures}
\begin{columns}
\begin{column}{0.6\textwidth}
\concept{Multiplexer Characteristics}:
\begin{itemize}
\item \textbf{Structure}: Tree of 2:1 multiplexers
\item \textbf{Delay}: Logarithmic with number of inputs
\item \textbf{Area}: Linear with number of inputs
\item \textbf{Control}: Binary encoded select signals
\end{itemize}

\vspace{0.3cm}
\concept{When to Use Multiplexers}:
\begin{itemize}
\item Mutually exclusive conditions
\item Fixed timing requirements
\item Parallel data selection
\item Performance-critical paths
\end{itemize}

\vspace{0.3cm}
\concept{Design Considerations}:
\begin{itemize}
\item Input data width affects area
\item Select signal encoding impacts delay
\item Technology-specific optimizations
\item Power consumption for unused inputs
\end{itemize}
\end{column}
\begin{column}{0.4\textwidth>
\begin{center}
\textbf{4:1 Multiplexer Tree}
\vspace{0.3cm}

\tikz{
% Inputs
\foreach \y/\label in {3.5/A, 3/B, 2.5/C, 2/D} {
  \node at (0,\y) {\label};
  \draw[->] (0.2,\y) -- (0.8,\y);
}

% First level muxes
\draw[thick] (1,3.5) rectangle (1.5,3);
\draw[thick] (1,2.5) rectangle (1.5,2);
\node at (1.25,3.25) {MUX};
\node at (1.25,2.25) {MUX};

% Second level mux
\draw[thick] (2,3) rectangle (2.5,2.5);
\node at (2.25,2.75) {MUX};

% Connections
\draw[->] (1.5,3.25) -- (2,2.9);
\draw[->] (1.5,2.25) -- (2,2.6);

% Select signals
\draw[->] (1.25,2.9) -- (1.25,2.7);
\node at (1.25,2.6) {S0};
\draw[->] (2.25,2.4) -- (2.25,2.2);
\node at (2.25,2.1) {S1};

% Output
\draw[->] (2.5,2.75) -- (3,2.75);
\node at (3.2,2.75) {Out};

% Timing
\node at (1.75,1.5) {Delay: $2 \times t_{mux}$};
}
\end{center}
\end{column}
\end{columns}
\end{frame}

\begin{frame}{Priority Logic Structures}
\begin{columns}
\begin{column}{0.6\textwidth}
\concept{Priority Logic Characteristics}:
\begin{itemize}
\item \textbf{Structure}: Cascaded conditional gates
\item \textbf{Delay}: Linear with condition position
\item \textbf{Area}: Depends on condition complexity
\item \textbf{Control}: Sequential evaluation order
\end{itemize}

\vspace{0.3cm}
\concept{When to Use Priority Logic}:
\begin{itemize}
\item Natural priority ordering
\item Overlapping conditions
\item Default/fallback cases
\item Early termination benefits
\end{itemize}

\vspace{0.3cm}
\concept{Optimization Opportunities}:
\begin{itemize}
\item Order by condition probability
\item Simplify condition expressions
\item Use early termination
\item Minimize critical path conditions
\end{itemize}
\end{column}
\begin{column}{0.4\textwidth>
\begin{center}
\textbf{Priority Logic Chain}
\vspace{0.3cm}

\tikz{
% Conditions
\foreach \y/\label in {3.5/C1, 3/C2, 2.5/C3, 2/Default} {
  \node at (0,\y) {\label};
}

% Logic gates
\draw[thick] (1,3.5) rectangle (1.8,3.2);
\node at (1.4,3.35) {AND};
\draw[thick] (1,3) rectangle (1.8,2.7);
\node at (1.4,2.85) {AND};
\draw[thick] (1,2.5) rectangle (1.8,2.2);
\node at (1.4,2.35) {AND};

% OR gate
\draw[thick] (2.5,3) rectangle (3.2,2.2);
\node at (2.85,2.6) {OR};

% Connections
\draw[->] (0.2,3.5) -- (1,3.4);
\draw[->] (0.2,3) -- (1,2.9);
\draw[->] (0.2,2.5) -- (1,2.4);

\draw[->] (1.8,3.35) -- (2.5,2.9);
\draw[->] (1.8,2.85) -- (2.5,2.7);
\draw[->] (1.8,2.35) -- (2.5,2.5);
\draw[->] (0.2,2) -- (2.5,2.3);

% Inhibit connections
\draw[->] (1.4,3.2) -- (1.4,3);
\draw[->] (1.4,2.7) -- (1.4,2.5);

% Output
\draw[->] (3.2,2.6) -- (3.7,2.6);
\node at (3.9,2.6) {Out};

% Timing
\node at (2,1.5) {Delay: Variable};
}
\end{center}
\end{column}
\end{columns}
\end{frame}

\begin{frame}{Decoder Logic Structures}
\begin{columns}
\begin{column}{0.6\textwidth}
\concept{Decoder Characteristics}:
\begin{itemize}
\item \textbf{Structure}: Binary to one-hot conversion
\item \textbf{Delay}: Fixed for all outputs
\item \textbf{Area}: Exponential with input width
\item \textbf{Control}: Distributed enable signals
\end{itemize}

\vspace{0.3cm}
\concept{When to Use Decoders}:
\begin{itemize}
\item Wide selector signals
\item One-hot control requirements
\item Distributed control logic
\item Memory address decoding
\end{itemize}

\vspace{0.3cm}
\concept{Implementation Considerations}:
\begin{itemize}
\item Input encoding affects complexity
\item Output loading impacts performance
\item Enable signals for power optimization
\item Technology-specific implementations
\end{itemize}
\end{column}
\begin{column}{0.4\textwidth>
\begin{center}
\textbf{2:4 Decoder}
\vspace{0.3cm}

\tikz{
% Inputs
\node at (0,3) {A1};
\node at (0,2.5) {A0};
\draw[->] (0.2,3) -- (1,3);
\draw[->] (0.2,2.5) -- (1,2.5);

% Decoder box
\draw[thick] (1,2) rectangle (2.5,3.5);
\node at (1.75,2.75) {2:4};
\node at (1.75,2.5) {DEC};

% Outputs
\foreach \y/\label in {3.3/Y3, 3/Y2, 2.7/Y1, 2.4/Y0} {
  \draw[->] (2.5,\y) -- (3.2,\y);
  \node at (3.4,\y) {\label};
}

% Truth table
\node at (1.75,1.5) {\textbf{One-hot outputs}};
\node at (1.75,1.2) {Only one active};

% Enable
\node at (0,1.8) {EN};
\draw[->] (0.2,1.8) -- (1,2.2);
}
\end{center}
\end{column}
\end{columns}
\end{frame}

% Section 5: Optimization Strategies
\section{Optimization Strategies}

\begin{frame}{Condition Ordering and Probability}
\begin{columns}
\begin{column}{0.6\textwidth}
\concept{Probability-Based Optimization}:
\begin{itemize}
\item Order conditions by execution frequency
\item Place most likely conditions first in \texttt{when}
\item Minimize average-case delay
\item Consider worst-case timing constraints
\end{itemize}

\vspace{0.3cm}
\concept{Condition Complexity Analysis}:
\begin{itemize}
\item Simple conditions before complex ones
\item Minimize logic depth for critical paths
\item Balance condition evaluation cost
\item Consider synthesis tool capabilities
\end{itemize}

\vspace{0.3cm}
\concept{Early Termination Benefits}:
\begin{itemize}
\item Reduce power consumption
\item Improve average-case performance
\item Simplify downstream logic
\item Enable aggressive optimizations
\end{itemize}
\end{column}
\begin{column}{0.4\textwidth>
\begin{center}
\textbf{Condition Ordering}
\vspace{0.3cm}

\tikz{
% Probability bars
\node at (1.5,3.5) {\textbf{Execution Frequency}};

\draw[thick, fill=green!30] (0,3) rectangle (2.5,3.2);
\node at (-0.5,3.1) {C1};
\node at (3,3.1) {80\%};

\draw[thick, fill=blue!30] (0,2.6) rectangle (1.5,2.8);
\node at (-0.5,2.7) {C2};
\node at (3,2.7) {15\%};

\draw[thick, fill=red!30] (0,2.2) rectangle (0.8,2.4);
\node at (-0.5,2.3) {C3};
\node at (3,2.3) {5\%};

% Optimization arrow
\draw[->, thick] (1.5,1.8) -- (1.5,1.4);
\node at (1.5,1.2) {\textbf{Order by frequency}};
\node at (1.5,0.9) {Most likely first};
}
\end{center}
\end{column}
\end{columns}
\end{frame}

\begin{frame}{Synthesis Tool Optimization}
\begin{columns}
\begin{column}{0.6\textwidth}
\concept{Tool-Specific Optimizations}:
\begin{itemize}
\item Leverage synthesis tool capabilities
\item Use tool-specific optimization directives
\item Understand tool limitations and preferences
\item Profile different coding styles
\end{itemize}

\vspace{0.3cm}
\concept{Optimization Passes}:
\begin{itemize}
\item Logic optimization and minimization
\item Technology mapping and binding
\item Timing-driven optimization
\item Power-aware transformations
\end{itemize}

\vspace{0.3cm}
\concept{Design Constraints}:
\begin{itemize}
\item Timing constraints and clock periods
\item Area and resource limitations
\item Power budgets and requirements
\item Technology-specific rules
\end{itemize}
\end{column}
\begin{column}{0.4\textwidth>
\begin{center}
\textbf{Optimization Flow}
\vspace{0.3cm}

\tikz{
% Optimization stages
\node[draw, rectangle, minimum width=2cm, minimum height=0.5cm] (logic) at (0,3) {Logic Opt};
\node[draw, rectangle, minimum width=2cm, minimum height=0.5cm] (tech) at (0,2.4) {Tech Map};
\node[draw, rectangle, minimum width=2cm, minimum height=0.5cm] (timing) at (0,1.8) {Timing Opt};
\node[draw, rectangle, minimum width=2cm, minimum height=0.5cm] (power) at (0,1.2) {Power Opt};

\draw[->] (logic) -- (tech);
\draw[->] (tech) -- (timing);
\draw[->] (timing) -- (power);

% Feedback loops
\draw[->, dashed] (timing.west) -- ++(-0.5,0) -- ++(0,0.6) -- (tech.west);
\draw[->, dashed] (power.west) -- ++(-0.7,0) -- ++(0,1.8) -- (logic.west);

% Constraints
\node[right] at (2.2,2.4) {Constraints};
\node[right] at (2.2,2.1) {Guide};
\node[right] at (2.2,1.8) {Optimization};
}
\end{center}
\end{column}
\end{columns}
\end{frame}

\begin{frame}{Best Practices for Performance Optimization}
\begin{columns}
\begin{column}{0.6\textwidth}
\concept{Design Guidelines}:
\begin{itemize}
\item Choose appropriate control construct for use case
\item Consider timing, area, and power trade-offs
\item Profile different implementation options
\item Use synthesis reports for optimization guidance
\end{itemize}

\vspace{0.3cm}
\concept{Coding Recommendations}:
\begin{itemize}
\item Use \texttt{MuxCase} for mutually exclusive conditions
\item Use \texttt{when} for priority-based logic
\item Order conditions by probability and complexity
\item Minimize condition evaluation complexity
\end{itemize}

\vspace{0.3cm}
\concept{Verification and Validation}:
\begin{itemize}
\item Verify functional equivalence across optimizations
\item Validate timing and power requirements
\item Test corner cases and edge conditions
\item Use formal verification when appropriate
\end{itemize}
\end{column}
\begin{column}{0.4\textwidth>
\begin{center}
\textbf{Optimization Checklist}
\vspace{0.3cm}

\tikz{
% Checklist items
\node at (1.5,3.5) {\textbf{Design Checklist}};

\foreach \y/\item in {3.2/Construct Choice, 2.9/Condition Order, 2.6/Timing Analysis, 2.3/Area Estimation, 2/Power Analysis, 1.7/Synthesis Reports, 1.4/Verification} {
  \draw[thick] (0,\y-0.1) rectangle (0.2,\y+0.1);
  \node[right] at (0.3,\y) {\item};
}

% Checkmarks
\foreach \y in {3.2, 2.9, 2.6, 2.3, 2, 1.7, 1.4} {
  \node at (0.1,\y) {✓};
}
}
\end{center}
\end{column}
\end{columns}
\end{frame}

% Section 6: Design Trade-offs
\section{Design Trade-offs}

\begin{frame}{Performance vs. Area Trade-offs}
\begin{columns}
\begin{column}{0.6\textwidth}
\concept{Trade-off Analysis}:
\begin{itemize}
\item \textbf{High Performance}: Parallel logic, larger area
\item \textbf{Low Area}: Sequential logic, longer delays
\item \textbf{Balanced}: Hybrid approaches, moderate both
\item \textbf{Application-Specific}: Optimize for critical metrics
\end{itemize}

\vspace{0.3cm}
\concept{Design Space Exploration}:
\begin{itemize}
\item Evaluate multiple implementation options
\item Use Pareto-optimal solutions
\item Consider application requirements
\item Balance competing objectives
\end{itemize}

\vspace{0.3cm}
\concept{Optimization Strategies}:
\begin{itemize}
\item Pipeline critical paths for performance
\item Share resources for area efficiency
\item Use clock gating for power savings
\item Leverage tool-specific optimizations
\end{itemize}
\end{column}
\begin{column}{0.4\textwidth>
\begin{center}
\textbf{Design Space}
\vspace{0.3cm}

\tikz{
% Axes
\draw[->] (0,0) -- (3,0);
\draw[->] (0,0) -- (0,3);
\node at (3.2,0) {Area};
\node at (0,3.2) {Performance};

% Pareto frontier
\draw[thick, blue] (0.5,2.8) .. controls (1.5,2) .. (2.8,0.5);

% Design points
\fill[red] (0.7,2.5) circle (0.05);
\node at (0.7,2.8) {MuxCase};
\fill[red] (2.2,1.2) circle (0.05);
\node at (2.2,0.9) {when};
\fill[red] (1.5,1.8) circle (0.05);
\node at (1.5,2.1) {switch};

% Optimal region
\draw[thick, green, dashed] (0.5,2.8) .. controls (1.5,2) .. (2.8,0.5);
\node at (2,2.5) {Pareto};
\node at (2,2.2) {Frontier};
}
\end{center}
\end{column}
\end{columns}
\end{frame}

% Summary
\section{Summary}

\begin{frame}{Chapter Summary}
\begin{columns}
\begin{column}{0.6\textwidth}
\concept{Key Concepts Covered}:
\begin{itemize}
\item Performance optimization principles in digital logic
\item Chisel control constructs and their synthesis implications
\item Hardware structures: multiplexers, priority logic, decoders
\item Timing, area, and power trade-offs
\item Optimization strategies and best practices
\item Tool-specific considerations and constraints
\end{itemize}

\vspace{0.3cm}
\concept{Design Guidelines}:
\begin{itemize}
\item Choose constructs based on application requirements
\item Consider synthesis implications early in design
\item Profile multiple implementation options
\item Use synthesis tools effectively
\item Validate optimization results
\end{itemize}

\vspace{0.3cm}
\concept{Future Considerations}:
\begin{itemize}
\item Emerging synthesis technologies
\item AI-assisted optimization
\item Power-efficient design techniques
\item Advanced verification methods
\end{itemize}
\end{column}
\begin{column}{0.4\textwidth>
\begin{center}
\textbf{Optimization Summary}
\vspace{0.3cm}

\tikz{
\node[draw, circle, minimum width=1.5cm] (center) at (0,0) {Optimized};
\node[draw, circle, minimum width=1.5cm] (center2) at (0,0) {Design};
\node[draw, rectangle] (timing) at (0,1.5) {Timing};
\node[draw, rectangle] (area) at (-1.5,0.5) {Area};
\node[draw, rectangle] (power) at (1.5,0.5) {Power};
\node[draw, rectangle] (tools) at (0,-1.5) {Tools};

\draw[->] (timing) -- (center);
\draw[->] (center) -- (area);
\draw[->] (center) -- (power);
\draw[->] (center) -- (tools);
}
\end{center}
\end{column}
\end{columns}
\end{frame}

\begin{frame}{Discussion Questions}
\begin{enumerate}
\item How do different Chisel constructs affect synthesis results?
\item What factors should guide the choice between \texttt{when} and \texttt{MuxCase}?
\item How can condition ordering optimize performance in priority logic?
\item What are the trade-offs between parallel and sequential logic structures?
\item How do synthesis tools influence optimization strategies?
\item What role does application profiling play in optimization decisions?
\item How can power consumption be minimized in control logic?
\item What verification challenges arise from optimization transformations?
\end{enumerate}
\end{frame}

